\paragraph{Articulación de las coordenadas regionales y la incidencia de clausura.}
\label{inc-ampl-articulacion}
La función de esta prolongación se comprende a partir de una distinción
entre tres operaciones: conservar un registro, evaluar una coordenada y
extraer una relación de incidencia. Las tres actúan sobre información
producida por APP--TRIT--TPK. Su articulación hace explícita la tesis
estructural del artículo: una realización numérica posee una procedencia
operatoria, y determinados datos de esa procedencia admiten, además, una
realización combinatoria y reticular. Las ecuaciones permiten seguir ambos
recorridos con sus dominios y con la información que permanece en cada uno.

El punto de partida es el estado terminal con sus prolongaciones
compatibles. En él se conservan los residuos y cocientes de las dos hojas
APP, el régimen y la orientación TRIT, y el transporte TPK con su memoria.
Las regiones de clausura, propagación y autoescala se determinan en ese
dominio antes de que sus lectores establezcan las coordenadas
$\pi,e,\varphi$. Al mismo tiempo, el registro orientado de las doce
ventanas determina el vector firmado $U$. La aparición posterior de $K$
procede de la transformación integral demostrada en
\S\ref{subsec:k-hadamard}:
\begin{equation}
 U=\mathcal H_{12}K,\qquad
 K=\frac14\mathcal H_{12}U,\qquad
 \mathcal H_{12}^{2}=4I_{12}.
 \label{inc-ampl-registro}
\end{equation}
La congruencia que define $\mathcal L_H$ garantiza la integralidad de
la inversa. De este modo, $K$ constituye una coordenatización reversible
del registro firmado en la subred pertinente. Sus posiciones conservan el
orden de fase; sus diferencias $D_3K,D_4K$ conservan las relaciones
transversales; la suma $Q(K)$ determina la componente uniforme. Esa
estructura explica por qué el registro interviene tanto en la composición
numérica como en la selección de una configuración de frontera.

La evaluación periódica
\[
 K\longmapsto\kappa_{\mathrm{per}}
 =\frac{\sum_{m=1}^{12}K_m1000^{12-m}}{1000^{12}-1}
\]
mantiene recuperables los doce bloques cuando se fijan base, longitud y
origen de fase. La extracción de un soporte realiza una operación
diferente: conserva la pertenencia de cada posición a una clase definida
por un umbral y agrupa los registros que tienen el mismo soporte. La
distinción entre estas dos aplicaciones resulta decisiva. La precisión
escalar de $\kappa_{\mathrm{per}}$ y la incidencia de $B_K$ son dos
contenidos matemáticos específicos, relacionados mediante el registro
completo del que proceden.

